% !TEX program = xelatex
% =============================================================================
% cn-everyday · 授权委托书（民事代办通用，律所级要件）
%
% 按律师起草标准写全：双方身份全要素（住址/电话）、法律依据、委托事项
% 逐项列明（对象特定化：写明不动产坐落与权证号）、权限边界（承认后果＋
% 明示排除事项——代收款项是最大风险点，样张明确排除）、转委托禁止、
% 期限（区间＋办结孰先）、生效要件与份数、附身份证件声明。
% 仅适用于一般民事代办；诉讼委托另有法院格式，勿混用。
%
% 编译：bash scripts/compile.sh templates/cn-everyday/power-of-attorney.tex --preview
% =============================================================================
\documentclass[zihao=-4]{ctexart}
\usepackage{everyday}

\begin{document}

\edTitle{授权委托书}

委托人：\hl{林××}，身份证号：\edblank[10em]{110×××××××××××××××}，
住址：示例市示例区示例路12号，联系电话：\edblank[8em]{138\,0000\,0000}。

受托人：\hl{张××}，身份证号：\edblank[10em]{110×××××××××××××××}，
住址：示例市示例区示例街33号，联系电话：\edblank[8em]{139\,0000\,0000}。

委托人因公务原因无法亲自到场办理下列事项，现依据《中华人民共和国
民法典》关于委托代理的规定，授权受托人在委托权限内代为办理：

一、代为办理委托人名下坐落于\hl{示例市示例区示例路100号}房屋
（不动产权证号：示例（2026）第000×××号）出售的网签备案、缴税及
不动产权属转移登记手续；

二、代为签署与上述事项相关的申请表、询问笔录、确认书及税费申报
文件；

三、代为领取完税凭证、不动产登记证明及相关票据文件。

\hl{权限与限制：}受托人在上述委托权限内实施的行为及签署的文件，
委托人均予承认，由此产生的法律后果由委托人承担；\hl{受托人无权代为
收取售房款项，无权变更交易价款，无转委托权。}

\hl{委托期限：}自2026年7月18日起至2026年10月17日止，或至上述事项
全部办结之日止（以先到者为准）。

本委托书自委托人签名并按手印之日起生效，一式三份，委托人、受托人
及登记机关各执一份。附：委托人、受托人身份证复印件各一份。

\vspace{0.8\baselineskip}
\edSign{委托人（签名并按手印）：}{林××}
\edSign[9em]{日期：}{2026年7月18日}

\end{document}
